% !TeX root = ../main.tex

\sysusetup{
  keywords  = {兴趣点推荐；生成式推荐；语义标识；地理信息增强；大语言模型},
  keywords* = {point-of-interest recommendation; generative recommendation; semantic ID; geographical enhancement; large language models},
}

\begin{abstract}
下一个兴趣点推荐（Next POI Recommendation）是基于位置服务和个性化推荐领域中的重要研究问题，对于地图搜索、出行服务和本地生活平台具有重要应用价值。近年来，随着大语言模型（Large Language Models, LLMs）和生成式推荐方法的发展，将推荐任务重构为序列生成问题逐渐成为一种新的研究范式。GNPR-SID 通过残差向量量化变分自编码器（Residual Quantized Variational AutoEncoder, RQ-VAE）学习兴趣点的语义标识（Semantic ID, SID），并利用大语言模型根据用户历史访问序列生成下一兴趣点的 SID，从而实现生成式 POI 推荐。然而，现有方法在 Semantic ID 构造阶段对地理信息的利用仍然较为有限，难以充分刻画兴趣点之间的空间差异及用户行为中的地理迁移规律。

针对上述问题，本文在 GNPR-SID 框架基础上提出了一种基于地理信息增强表示的生成式 Next POI 推荐方法。该方法不再直接使用原始的类别名称和区域信息构造兴趣点连续表示，而是首先利用轻量文本编码器 \texttt{all-MiniLM-L6-v2} 对兴趣点类别名称进行编码，获得基础语义向量；然后结合经纬度信息，通过参考点聚类、方位角计算和旋转位置编码构造地理增强表示（geo\_emb）；最后将 geo\_emb 输入 RQ-VAE 学习具有空间感知能力的 SID，并在后续生成式推荐阶段利用大语言模型根据历史 SID 序列预测下一兴趣点。与显式地将地理编码作为额外输入或输出目标的方式不同，本文方法强调在表示学习阶段隐式融入地理信息，从而在保持生成框架简洁性的同时提升兴趣点离散表示质量。

本文在 NYC、TKY 和 CA 三个数据集上进行了实验验证，并与基线方法进行了对比分析。实验结果表明，在本文的实验环境与复现设置下，本文方法在三个数据集上的 Acc@1 分别达到 0.3550、0.2829 和 0.2122，均较对应的基线模型取得提升，说明在 Semantic ID Construction 阶段引入地理增强表示能够有效改善生成式 Next POI 推荐性能。进一步的消融实验表明，地理信息本身对 POI 推荐任务具有重要作用，而相比显式地理信息注入方式，本文采用的隐式地理增强方法在预测效果和建模效率之间取得了更好的平衡。研究结果表明，在生成式 POI 推荐场景下，通过改进离散语义表示学习过程引入空间信息，是提升模型空间感知能力和推荐性能的一种有效途径。

\end{abstract}

\begin{abstract*}
Next point-of-interest recommendation is an important research problem in location-based services and personalized recommendation, and it is of great practical value for map search, travel assistance, and local lifestyle platforms. In recent years, with the rapid development of large language models and generative recommendation methods, reformulating recommendation as a sequence generation task has become a promising paradigm. GNPR-SID learns semantic IDs for POIs through a Residual Quantized Variational AutoEncoder (RQ-VAE), and then fine-tunes a large language model to generate the semantic ID of the next POI based on a user’s historical check-in sequence. However, existing methods still make limited use of geographical information during semantic ID construction, making it difficult to sufficiently capture spatial differences among POIs and geographical transition patterns in user mobility.

To address this issue, this thesis proposes a geography-enhanced generative next POI recommendation method based on the GNPR-SID framework. Instead of directly using raw category names and region information to construct POI representations, the proposed method first employs the lightweight text encoder \texttt{all-MiniLM-L6-v2} to encode POI category names into base semantic vectors. Then, by incorporating latitude and longitude information, a geography-enhanced embedding (geo\_emb) is constructed through reference point clustering, bearing angle computation, and rotational positional encoding. Finally, geo\_emb is fed into RQ-VAE to learn geography-aware semantic IDs, which are further used by a large language model to predict the next POI from historical SID sequences. Different from approaches that explicitly inject geographical codes into input or output sequences, the proposed method implicitly integrates spatial information during the representation learning stage, improving POI discrete representation quality while keeping the generative framework concise.

Experiments are conducted on three datasets, namely NYC, TKY, and CA. The results show that, under the experimental environment and reproduction settings adopted in this thesis, the proposed method achieves Acc@1 scores of 0.3550, 0.2829, and 0.2122 on the three datasets, respectively, consistently improving over the corresponding baseline results. This indicates that introducing geography-enhanced representations into the Semantic ID Construction stage can effectively improve generative next POI recommendation. Further ablation studies demonstrate that geographical information plays an important role in POI recommendation, and that the proposed implicit geography-enhanced strategy provides a better balance between predictive performance and modeling efficiency than explicit geographical injection. The findings suggest that, in generative POI recommendation, introducing spatial information through improved discrete semantic representation learning is an effective way to enhance both spatial awareness and recommendation performance.

\end{abstract*}